% % % % % % % % % % % % % % % % % % % % % % % % % % % % % % % % % % % % % % % % % % % % % % % % % % % % % % % % % % % % % % % %
% SAT_STYLE_MATH_MCQS.tex
% 1_14.tex
% 
% encoding: UTF-8 
% EOL: LF
%
% format: LaTeX
% indent: spaces (2)
% width: 127
% % % % % % % % % % % % % % % % % % % % % % % % % % % % % % % % % % % % % % % % % % % % % % % % % % % % % % % % % % % % % % % %
Une montre à aiguilles indique $16$\,h\,$16$: quel est l'angle formé par
l'aiguille des heures et celle des minutes?
\begin{enumerate}
  \item $30°$
  \item $32°$
  \item $34°$
  \item $36°$
  \item $38°$
\end{enumerate}
%
%
\begin{proof}[\textbf{B}]
À $16$\,h\,$00$, l'angle est de $120°$. En $16$ minutes:
\begin{itemize}
  \item l'aiguille des heures avance de $\frac{16 \times 30°}{60} = 8°$;
  \item l'aiguille des minutes avance de $8° \times 12 = 96°$.
\end{itemize}
L'angle vaut donc $120° + 8° - 96° = 32°$.
\end{proof}